The exposure entering an allowance calculation also needs a defined
population. JPMorgan's 2025 report gives four different amounts for
lending-related commitments at 31 December 2025:
\begin{center}
\begin{tabularx}{\linewidth}{@{}Xr@{}}
\toprule
Reported quantity & US\$ millions \\
\midrule
Contractual lending-related amount & 1,817,307 \\
Commitments in the allowance-methodology table & 603,947 \\
Allowance for lending-related commitments & 5,071 \\
Lending-related carrying value & 6,278 \\
\bottomrule
\end{tabularx}
\end{center}
Notes 13 and 28 define these populations and measurements separately
\citep[pp.~258--262, 295--297]{jpmorgan2025tenk}. The contractual total
contains consumer and wholesale commitments; the allowance table excludes
specified populations. The commitment allowance is reported in accounts
payable and other liabilities. Carrying value can also contain guarantee
liabilities, fair-value measurements and deferred fees or purchase adjustments,
as applicable. For standby letters of credit and other financial guarantees,
the report explicitly reconciles a US\$175m allowance and a US\$441m guarantee
liability to US\$616m of carrying value. Treating either the contractual total
or carrying value as the allowance would compute the wrong quantity.

The report also distinguishes contractual capacity from expected funding.
Customers may leave facilities unused, refinance them or let them expire;
some commitments can be reduced or cancelled under their contracts or law.
The reported policy records no allowance for future draws on unconditionally
cancellable commitments. That exclusion concerns the loss measurement;
a liquidity scenario still needs the conditions under which a customer can
draw, the bank can cancel, and a payment requires settlement funds. Those
conditions and their dependence on borrower distress cannot be recovered
from a commitment total. Comparisons also need the same reporting basis:
specified wholesale amounts are net of risk participations here but gross
in Federal Reserve filings. Parent guarantees of subsidiary debt are
expressly outside the Note~28 table
\citep[pp.~258, 295--296, 300]{jpmorgan2025tenk}.

A commitment to acquire loans makes the time distinction especially clear.
The report dates the Apple Card forward purchase commitment to 30 December
2025, its announcement to 7 January 2026, and expected closing to
approximately 24 months later. The commitment disclosure includes about
US\$104bn of estimated credit exposure at closing, including about US\$23bn
of estimated drawn loans, while the reported transaction allowance is
US\$2.2bn \citep[pp.~47, 260, 296]{jpmorgan2025tenk}. The expected drawn
loans at that future closing are not loans already owned at the 2025 reporting
date. A reconstruction therefore retains the forward commitment, the future
portfolio estimate and the dated allowance as different facts. A forecast
must supply closing and subsequent draw assumptions before generating loan
and settlement events. These reported policies and amounts support the
distinctions; their application to a new calculation still requires the
operative accounting rules and actual contract facts.

